\clearpage
\section{Appendix: Four One-Page Dogma Synopses}

These sheets are retrieval aids, not substitutes for the preceding argument or the cited primary texts. Each preserves the difference between dogmatic object, evidence, theological explanation, historical development, and an open question.

\subsection*{I. Divine Motherhood: Mary is truly \emph{Theotokos}}
\phantomsection
\addcontentsline{toc}{subsection}{I. Divine Motherhood: one-page synopsis}

\begingroup
\footnotesize
\renewcommand{\arraystretch}{1.12}
\begin{tabularx}{\textwidth}{@{}>{\raggedright\arraybackslash\bfseries}p{0.205\textwidth}X@{}}
\toprule
Exact object & Mary is truly Mother of God because the one conceived and born from her according to his humanity is the eternal Word: one divine person, truly God and truly man.\\
Mode and authority & Ecumenical conciliar Christological judgment: Ephesus (431) received Cyril's Second Letter and first anathema; the Formula of Union (433) and Chalcedon (451) give the balanced grammar of one person in two natures.\\
Scriptural center & Matt 1:18--25; Luke 1:31--35, 43; Gal 4:4; John 1:1--14. Scripture supplies the reality rather than the later Greek title: God's pre-existent Son is truly born of a woman and Elizabeth names Mary mother of her Lord.\\
Patristic support & Ignatius joins real virginal birth to real saving flesh; Irenaeus teaches recapitulation and the New Eve; Gregory Nazianzen explicitly defends \emph{Theotokos}; Cyril explains that the Word personally unites flesh from Mary.\\
Thomistic support & ST III, q.35, a.4: maternity is said with respect to the person, and the person conceived and born is divine. ST III, q.2 guards the hypostatic union without confusing the natures.\\
Historical hinge & Anti-docetic realism matures into fourth-century title and fifth-century conciliar precision. Ephesus does not make Mary divine; it excludes a divided Christ. Chalcedon's “according to the humanity” protects both the title and the distinction of natures.\\
Doctrinal fruit & The dogma secures the reality of the Incarnation, the communicatio idiomatum, and the saving value of the Son's human obedience and suffering. It is the Christological root of the other Marian privileges.\\
Assent and boundary & Received with divine and Catholic faith. Mary does not cause the divine nature, precede the Trinity, or become a goddess. “Mother of Christ” is adequate only when Christ unmistakably names the one divine person.\\
\bottomrule
\end{tabularx}

\begin{studybox}{Memory formula}
Mary's maternity reaches a person; the person born from her is God the Word incarnate, according to his humanity.
\end{studybox}
\endgroup
\clearpage
\subsection*{II. Perpetual Virginity: before, in, and after the birth}
\phantomsection
\addcontentsline{toc}{subsection}{II. Perpetual Virginity: one-page synopsis}

\begingroup
\footnotesize
\renewcommand{\arraystretch}{1.12}
\begin{tabularx}{\textwidth}{@{}>{\raggedright\arraybackslash\bfseries}p{0.205\textwidth}X@{}}
\toprule
Exact object & Mary conceived Christ without a human father, brought him forth with her virginal integrity preserved, and remained virgin thereafter: \emph{ante partum, in partu, post partum}.\\
Mode and authority & Dogma of universal Tradition and the ordinary universal Magisterium, repeatedly confessed by councils and popes. Constantinople II (553) calls her ever-Virgin; Lateran 649, canon 3, states the threefold content explicitly; LG 57 and CCC 496--507 synthesize it.\\
Scriptural center & Matt 1:18--25 and Luke 1:26--38 directly govern the virginal conception. Luke 2:7 affirms real birth. The Church reads “brothers,” “until,” and “firstborn” within biblical idiom and the received faith, not as proof of later children.\\
Patristic support & Ignatius witnesses virginal birth; Epiphanius uses \emph{Aeiparthenos}; Jerome answers Helvidius's exegesis; Augustine joins prior purpose, real maternity, and enduring virginity; Leo the Great confesses true birth with preserved integrity.\\
Thomistic support & ST III, q.28 treats virginity before, in, and after birth; q.32, a.3 denies paternity to the Holy Spirit while affirming conception by his power; q.35 preserves genuine motherhood and birth.\\
Historical hinge & The substance precedes the later three-part formula. Controversies sharpen the meaning; Lateran 649 supplies an explicit landmark, while later councils and Paul IV continue the confession. The dogma is not inferred from one isolated proof-text.\\
Doctrinal fruit & The sign reveals Christ's divine initiative, true Incarnation, and Mary's undivided consecration. It honors marriage and vowed virginity according to distinct vocations; it does not portray ordinary childbirth as defiling.\\
Assent and boundary & The Church does not define microscopic physiology, one exclusive genealogy for every “brother,” or a notarized canonical vow in Luke 1:34. Christ's birth is real, not a passage through an unreal body.\\
\bottomrule
\end{tabularx}

\begin{studybox}{Memory formula}
One undivided confession: conceived virginally, truly born without loss of virginal integrity, and no later marital consummation or children.
\end{studybox}
\endgroup

\clearpage
\subsection*{III. Immaculate Conception: preservative redemption}
\phantomsection
\addcontentsline{toc}{subsection}{III. Immaculate Conception: one-page synopsis}

\begingroup
\footnotesize
\renewcommand{\arraystretch}{1.10}
\begin{tabularx}{\textwidth}{@{}>{\raggedright\arraybackslash\bfseries}p{0.205\textwidth}X@{}}
\toprule
Exact object & From the first instant of her conception Mary was, by a singular grace and privilege of Almighty God and in view of Christ's merits, preserved immune from every stain of original sin.\\
Mode and authority & Ex cathedra definition by Pius IX, \emph{Ineffabilis Deus}, 8 December 1854. The constitution defines the doctrine as revealed by God and therefore to be firmly and constantly believed.\\
Scriptural center & Gen 3:15, Luke 1:28, the New Eve pattern, and the biblical economy of grace supply a canonical matrix; no verse states the 1854 formula word for word. Romans 5 remains controlling: Mary needs the one Savior and is redeemed more perfectly by preservation.\\
Patristic support & Justin and Irenaeus establish New Eve obedience; Ephrem and liturgical language witness singular holiness; Augustine excludes Mary when discussing personal sin. These are genuine antecedents, not verbatim definitions of first-instant preservation.\\
Thomistic limit and use & ST III, q.27 and \emph{Compendium} I, c.224 place sanctification after animation and therefore do not teach the defined object. Thomas nevertheless guards Christ's universal mediation, Mary's fullness of grace, freedom, and lifelong holiness---principles preserved by the final formula.\\
Development & Conception feast; Eadmer and Bernard; scholastic dispute; Scotus's preservative-redemption logic; Sixtus IV; Trent; Paul V and Gregory XV; Alexander VII; Pius IX's 1849 episcopal consultation; definition in 1854. Disciplinary restrictions before 1854 are not retroactive definitions.\\
Doctrinal fruit & Mary is the first redeemed, wholly dependent upon the Paschal mystery. Grace perfects rather than replaces her freedom; her fiat is both grace-enabled and genuinely hers. She images the holiness Christ intends for the Church.\\
Assent and boundary & Not Christ's virginal conception, not Mary's conception without marital intercourse, not independence from Christ, and not identical to the separate doctrine of lifelong personal sinlessness. Lourdes can receive the dogma devotionally but cannot establish it.\\
\bottomrule
\end{tabularx}

\begin{studybox}{Memory formula}
The one Savior redeems Mary not after contraction of original sin but by preventing it from the first instant through his foreseen merits.
\end{studybox}
\endgroup

\clearpage
\subsection*{IV. Assumption: body and soul in heavenly glory}
\phantomsection
\addcontentsline{toc}{subsection}{IV. Assumption: one-page synopsis}

\begingroup
\footnotesize
\renewcommand{\arraystretch}{1.10}
\begin{tabularx}{\textwidth}{@{}>{\raggedright\arraybackslash\bfseries}p{0.205\textwidth}X@{}}
\toprule
Exact object & The Immaculate Mother of God and ever-Virgin Mary, having completed the course of her earthly life, was assumed body and soul into heavenly glory.\\
Mode and authority & Ex cathedra definition by Pius XII, \emph{Munificentissimus Deus} 44, 1 November 1950, following the 1946 episcopal consultation \emph{Deiparae Virginis Mariae}.\\
Scriptural center & 1 Cor 15:20--57, Rom 8:11, 23, and Phil 3:20--21 directly reveal bodily resurrection and glory in Christ. Luke 1, the New Eve, Ark imagery, Ps 132:8, and Rev 12 provide Marian fittingness or typology, not a hidden canonical Assumption narrative.\\
Patristic and liturgical support & Epiphanius honestly records early uncertainty about Mary's end. Later \emph{Transitus} texts show reception but are not eyewitness memoirs. Dormition/Assumption worship and the mature homilies of Pseudo-Modestus, Germanus, and John Damascene clearly confess bodily glorification.\\
Thomistic support and limit & Aquinas has no dedicated \emph{Summa} proof; Pius XII cites occasional affirmative loci. Thomistic anthropology, the soul as substantial form, and Christ's Resurrection as cause and exemplar illuminate why salvation culminates bodily without serving as historical proof.\\
Development & Late-antique narrative traditions; Eastern Dormition and Roman feast; medieval consensus; increasing papal teaching; worldwide episcopal consultation; the 1950 definition. Development does not canonize every tomb, house, date, homily, or apocryphal detail.\\
Doctrinal fruit & The Assumption is an anticipated participation in Christ's Easter victory and an icon of the Church's destiny. It honors embodied creaturehood and transformed communion; it does not create a second resurrection-source beside Christ.\\
Assent and boundary & “Was assumed” names God's action on a redeemed creature, unlike Christ's Ascension by his own power. The definition leaves Mary's death open; Dormition remains a weighty tradition. It does not define mechanism, date, place, tomb, or omnipresence.\\
\bottomrule
\end{tabularx}

\begin{studybox}{Memory formula}
Christ is the risen firstfruits and cause; Mary is the redeemed creature whose whole embodied person already shares his promised glory.
\end{studybox}
\endgroup
\clearpage
